% page 364 of the facsimile (image p-363 of the scan)
% block 154.9 x 229.0 mm ; left margin 8.0 mm
\pgc{-12.700mm}{0.000mm}{
\l{\cel{3}356.}
\l{}
\l{}
\l{\sou{mapo.}\cel{1}Folio ek papero rezistiva, ma flexebla (facita ek}
\l{\cel{3}plura paperfolii kunglutinita) od ek papero, sur qua}
\l{\cel{3}esas reprezentita la surfaco di la Terglobo o di un ek}
\l{\cel{3}olua parti. - E.}
\l{}
\l{\sou{maro}. I. Extensajo vasta ek aquo saloza, qua okupas parto}
\l{\cel{3}grandega di la surfaco di la Terglobo. - II. Parto di}
\l{\cel{3}ta extensajo vasta ek aquo saloza, situita en ta o ca}
\l{\cel{3}regiono. - III. (\sou{metaf}.)Extensajo vasta di liquido -}
\l{\cel{3}ek aquo konjelita. (Uzata anke pri granda quanto de su-}
\l{\cel{3}cii, de desfacilaji en qui onu esas quaze imersita). -}
\l{\cel{3}DEFIRS.}
\l{}
\l{\sou{marabuto}. (\sou{religio di Islam}). Mohamedisto konsakrita a la}
\l{\cel{3}praktiko ed a la doco di sua religio. - DEFIRS.}
\l{}
\l{\sou{maraskino}. Ratafio qua fabrikesas en Italia, ek cerizi}
\l{\cel{3}acida. - DEFIRS.}
\l{}
\l{\sou{marasmo}. (\sou{patol}.)Langoro qua akompanas la deperiso, la sen-}
\l{\cel{3}kurajeso. - DEFIRS.}
\l{}
\l{\sou{marchar}. (\sou{netrans.}) Pozar un ek sua pedi, e pose la altra,}
\l{\cel{3}e tale ri-agante lo, por irar ta-direcione o ca-direcione.}
\l{\cel{3}- DEFIRS.}
\l{}
\l{\sou{marchandar}. (\sou{trans.)}\cel{2}Probar obtenor po preco min granda.-}
\l{\cel{3}F.}
\l{}
\l{\sou{marcipano}.\cel{2}Kukajo ek mandeli pistita e sukro. - DEFIRS.}
\l{}
\l{\sou{mardio}. Dio triesma di la semano, qua sequas lundio. - FIS.}
\l{}
\l{\sou{mareo}. Movo alternanta (di depreseso e di plualteso) di}
\l{\cel{3}maro qua, en 24 hori, retroiras pokope ed abandonas la}
\l{\cel{3}litoro, e pose ri-venas kovrar la spaco quan lu lasabis}
\l{\cel{3}sikeskar. - FIS.}
\l{}
\l{\sou{mareografo}. (\sou{fiziko)}. Instrumento por mezurar ed enrejistrar}
\l{\cel{3}la varii di la marei, per flotacero, trasante kurvo.qua}
\l{\cel{3}reprezentas la ampleso-grado di la mareo, sur paperbendo}
\l{\cel{3}senfina quan desvolvas mekanismo horlojala. - EF.}
\l{}
\l{\sou{mareometro}. Sinonimo di \textquotedbl{}mareografo\textquotedbl{}.}
\l{}
\l{\sou{margarino}. (\sou{kemio}). Kombinuro ek margar-acido e glicerino,}
\l{\cel{3}uzata kom butro. - DEFIS.}
\l{}
\l{\sou{margravo}. (\sou{historio}) Nomo di ula (olima) princi suverena}
\l{\cel{3}di Germania. - DEFIS.}
}
